\chapter{Expérimentation, résultats et discussion}
\label{chap:experimentations}

Le chapitre~2 a défini un protocole expérimental complet --- corpus, pipeline,
métriques d'efficacité et de frugalité --- sans toutefois formuler d'hypothèses
de recherche formelles, faute de mesures permettant de fixer les seuils de
dégradation acceptable sur une base empirique plutôt qu'arbitraire
(section~2.5.2). Ce chapitre met ce protocole à l'épreuve : il rapporte les
mesures pilotes qui permettent de formuler ces hypothèses, présente les résultats
obtenus pour chacun des modules du pipeline formalisé en section~2.4.1, puis en
propose une discussion critique.

Il poursuit trois objectifs successifs. D'abord, un objectif de calibrage :
conduire les mesures pilotes nécessaires à la formulation d'hypothèses de
recherche précises, plutôt que génériques ou arbitrairement fixées. Ensuite, un
objectif de mesure : rapporter, module par module --- reconnaissance manuscrite,
extraction sémantique, résolution d'entités et peuplement du graphe ---, la
qualité et la frugalité effectivement obtenues, présentées conjointement pour
chaque module conformément au cadre de dominance frugale posé en section~2.1.2.
Enfin, un objectif d'interprétation : replacer ces résultats chiffrés dans la
perspective critique attendue d'un travail de recherche, en particulier au regard
des lacunes identifiées au chapitre~1 et du compromis précision/frugalité au
cœur de la contribution de ce mémoire.

Le périmètre de ce chapitre appelle une précision. La discussion critique des
résultats est volontairement concentrée dans une section dédiée (section~3.6),
plutôt que dispersée au fil de leur présentation : cette dernière (sections~3.3
à~3.5) reste descriptive, au plus près des chiffres obtenus, pour que
l'interprétation --- nécessairement plus sujette à discussion --- soit clairement
identifiable comme telle plutôt que mêlée à la présentation des faits. La portée
générale de la contribution du mémoire, au-delà des résultats propres à ce
chapitre, ainsi que les perspectives de travaux futurs, ne sont pas traitées ici
mais dans la Conclusion générale qui suit ce chapitre.

Le chapitre est organisé en six sections. La section~3.1, déjà présentée, a établi
le cadre expérimental et les conditions de reproductibilité. La section~3.2
rapporte les mesures pilotes et formule les hypothèses de recherche qui en
résultent. Les sections~3.3, 3.4 et~3.5 présentent respectivement les résultats
du module de reconnaissance manuscrite, du module d'extraction sémantique, et de
la résolution d'entités et du peuplement du graphe de connaissances. La
section~3.6, enfin, discute ces résultats à la lumière des hypothèses formulées
en section~3.2 et du positionnement établi au chapitre~1.

\section{Cadre expérimental et reproductibilité}

\subsection{Environnement matériel et logiciel}

Les expérimentations sont conduites sur la même machine que celle mobilisée pour
la collecte du corpus (section~2.3.1) : un ThinkPad~T460, 12~Go de RAM, 256~Go de
stockage SSD, sans GPU dédié. Ce choix n'est pas fortuit : il incarne concrètement,
au niveau du cadre expérimental lui-même, le principe retenu en section~2.1.3 ---
une infrastructure ancienne et largement répandue plutôt qu'un équipement
spécialisé, cohérente avec la contrainte \og machine standard sans GPU \fg{} déjà
posée pour les métriques de frugalité (section~2.5.2).

% TODO : compléter le modèle exact de processeur et les versions logicielles
\begin{table}[H]
    \centering
    \small
    \caption{Environnement logiciel}
    \label{tab:environnement-logiciel}
    \begin{tabular}{ll}
        \toprule
        \textbf{Composant}                                      & \textbf{Version} \\
        \midrule
        Python                                                  & [à compléter]    \\
        Kraken                                                  & [à compléter]    \\
        \texttt{transformers} / \texttt{peft} (NER, LoRA/QLoRA) & [à compléter]    \\
        OpenCV (prétraitement, section~2.3.2)                   & [à compléter]    \\
        GraphDB (section~2.4.6)                                 & [à compléter]    \\
        \texttt{jiwer} (calcul CER/WER)                         & [à compléter]    \\
        \texttt{seqeval} (calcul P/R/$F_1$ NER)                 & [à compléter]    \\
        \bottomrule
    \end{tabular}
\end{table}

\subsection{Corpus de test et répartition entraînement/validation/test}

Le corpus mobilisé pour les expérimentations est celui décrit en section~2.3.1 :
1568~actes collectés à la mairie de Koumra (Tchad), majoritairement des actes de
naissance. Deux sous-ensembles annotés en sont extraits, selon les protocoles de
la section~2.3.3 : une vérité de terrain HTR via eScriptorium, et un jeu d'entités
et de relations annoté via Label Studio --- ces deux sous-ensembles ne portant pas
nécessairement sur les mêmes actes, l'annotation HTR opérant au niveau de la ligne
et l'annotation NER/RE au niveau de l'acte entier.

Pour le module HTR, la répartition entraînement/validation/test s'effectue au
niveau de la \emph{ligne}, unité de traitement native de Kraken (section~2.4.2),
selon une proportion de 70~\%/15~\%/15~\%, sous la contrainte déjà posée en
section~2.4.2 : les lignes d'un même acte ne peuvent apparaître dans plus d'un
ensemble à la fois, afin d'éviter toute fuite de données entre ensembles. Pour le
module NER/RE, la répartition s'effectue au niveau de l'\emph{acte}, dans la même
proportion, stratifiée par type d'acte pour préserver, dans chacun des trois
ensembles, la proportion relative de naissances, mariages et décès observée dans
le corpus global --- une stratification qui atténue, sans l'éliminer, le
déséquilibre entre types d'actes déjà relevé comme limite en section~2.5.3.

% TODO : compléter une fois l'annotation stabilisée
\begin{table}[H]
    \centering
    \small
    \caption{Répartition des sous-ensembles annotés}
    \label{tab:repartition-ensembles}
    \begin{tabular}{lccc}
        \toprule
        \textbf{Sous-ensemble}      & \textbf{Entraînement} & \textbf{Validation} & \textbf{Test} \\
        \midrule
        Lignes HTR (eScriptorium)   & [à compléter]         & [à compléter]       & [à compléter] \\
        Actes NER/RE (Label Studio) & [à compléter]         & [à compléter]       & [à compléter] \\
        \bottomrule
    \end{tabular}
\end{table}

\subsection{Protocole et outils de mesure}

Les métriques d'efficacité définies en section~2.5.1 sont calculées à l'aide
d'outils standard plutôt que d'implémentations propres, pour limiter le risque
d'erreur de calcul et faciliter la reproduction par un tiers. Le CER et le WER
sont calculés avec \texttt{jiwer}, sur la définition posée en section~2.5.1
(distance d'édition entre transcription et référence). La précision, le rappel et
le score $F_1$ du module NER sont calculés avec \texttt{seqeval}, au niveau de
l'entité (span et type identiques à la référence), plutôt qu'au niveau du token,
conformément à la section~2.5.1. Aucun outil standard n'existant pour l'évaluation
de relations définies sur mesure, celle du module RE repose sur un script
d'évaluation dédié, comparant les triplets (type, source, cible) produits par les
règles positionnelles (section~2.4.3) aux relations attendues, annotées manuellement
sur l'échantillon prévu à cet effet (section~2.5.1).

Les métriques de frugalité (section~2.5.2) suivent un protocole de mesure explicite
plutôt qu'une mesure ponctuelle isolée. Le temps de traitement par acte est mesuré
en excluant le chargement initial du modèle --- amorti une seule fois, non
représentatif du coût récurrent d'exploitation --- et rapporté sur plusieurs
exécutions répétées (médiane, minimum et maximum) plutôt que sur une seule mesure,
pour limiter l'effet du bruit de mesure inhérent à une machine partagée avec
d'autres processus. La mémoire vive est mesurée par l'empreinte mémoire résidente
maximale (RSS) du processus pendant l'inférence ; la mémoire vidéo, lorsqu'un GPU
est mobilisé --- exclusivement pour l'entraînement, non pour les mesures de
frugalité en inférence qui restent conduites sans GPU (section~2.5.2) ---, est
rapportée séparément et n'est pas comparable aux mesures de RAM. La taille sur
disque correspond à la taille du fichier de poids sauvegardé --- modèle complet
pour Kraken (section~2.4.2), adaptateur seul pour les configurations LoRA/QLoRA
du module NER (section~2.4.4), ce qui suppose de préciser, au moment de la
présentation des résultats (section~3.4.3), si la comparaison porte sur la taille
de l'adaptateur seul ou sur l'ensemble modèle de base~+~adaptateur nécessaire à
l'inférence.